\documentclass[12pt]{article}

\usepackage[a4paper, margin=1in]{geometry}
\usepackage{amsmath, amssymb}
\usepackage{setspace}
\usepackage{hyperref}

\setstretch{1.2}

\title{\textbf{Orch-OR Free Will Model: \\ A Computational Analogue of Quantum Stochastic Collapse in Neural Architectures}}
\author{Iyer Ramkumar Ramasubramanian \\ \textit{Independent Researcher}}
\date{\today}

\begin{document}

\maketitle

\begin{abstract}
The nature of free will remains a central problem at the intersection of physics and neurobiology. This paper introduces a computational analogue of the Orchestrated Objective Reduction (Orch-OR) theory proposed by Penrose and Hameroff. By integrating classical cortical processing with a "quantum-like" stochastic collapse mechanism and a recursive self-model, we demonstrate how non-deterministic yet structured agency can emerge within a neural substrate. Our model suggests that free will is effectively a form of structured stochasticity stabilized by a continuity of identity.
\end{abstract}

\section{Introduction}
As an independent researcher exploring the intersection of neurobiology and statistical mechanics, I propose that agency cannot be fully explained by deterministic classical physics. The \textbf{Orch-OR Free Will Model} extends the HDBM framework by introducing discrete "moments of choice" through objective reduction, moving beyond passive equilibrium into active, non-linear decision-making.

\section{Theoretical Core and Architecture}

The model simulates the Orch-OR process through a four-stage pipeline:

\begin{itemize}
\item \textbf{Classical Processing (Cortex)}: Traditional neural network layers handle feature extraction and environmental sensing.
\item \textbf{Quantum Collapse (Microtubules)}: A stochastic layer simulates the discrete reduction of possibilities into a single state, introducing non-deterministic variance into the system.
\item \textbf{Recursive Self-Model (DMN)}: A feedback loop analogous to the Default Mode Network ensures that each "collapse" is integrated into a persistent identity.
\item \textbf{Decision Output}: The final action is a "resolved" state that reflects both external data and internal spontaneity.
\end{itemize}

\section{System Dynamics}

The model distinguishes between simple randomness and true agency:

\subsection{Structured Stochasticity}
Unlike purely random models, the "collapse" in this system is constrained by the weights of the classical neural substrate. The system exhibits structured behavior that varies even when the initial state is identical.

\subsection{Moments of Choice}
Each time-step represents a discrete "moment of choice." The self-referential loop (Self $\rightarrow$ Decision $\rightarrow$ Self) prevents the system from diverging into chaos, maintaining a continuity of identity over time.

\section{Mathematical Formulation}

The state $S$ at time $t$ is no longer a deterministic function of inputs. Instead, it incorporates a collapse operator $\Psi$:

\[
S_{t+1} = \Psi(f(W S_t + X_t) + \epsilon_{q})
\]

Where $\epsilon_{q}$ represents the quantum-analogue variance. The transition is stabilized by the DMN-like recursive function $R$:

\[
Action = \text{Softmax}(R(S_{t+1}))
\]

\section{Inference and Observations}

Computational results of the Orch-OR model indicate:
\begin{itemize}
\item \textbf{Non-Deterministic Agency}: The agent chooses different actions in identical environments, proving the presence of non-deterministic choice.
\item \textbf{Path Stability}: Despite the stochasticity, the system does not collapse into zero action or random noise; it maintains a logical trajectory toward defined goals.
\item \textbf{Emergent Free Will}: Agency is shown to be a "resolved negotiation" between the laws of physics (weights) and the spontaneity of the collapse (stochasticity).
\end{itemize}

\section{Conclusion}
The Orch-OR Free Will Model provides a computational framework for investigating the biological basis of agency. By modeling the interaction between classical and quantum-like regimes, we offer a path toward understanding how the brain constructs a stable "self" within a non-deterministic universe.

\section{Repository for Experimentation}
\noindent
\url{https://github.com/spacetracker-collab/FreeWill}

\end{document}
